\section{Eigenvalue Problem of $\hat{L}^2$ and $\hat{L}_z$}

Having established that $[\hat{L}^2, \hat{L}_z] = 0$, we know there exists a common basis of eigenstates for these operators. We define these eigenstates as functions of the spherical angles, denoted by $Y(\theta, \phi)$.

Our goal is to solve the simultaneous eigenvalue equations:
\begin{align}
    \hat{L}_z Y(\theta, \phi) &= \lambda Y(\theta, \phi) \\
    \hat{L}^2 Y(\theta, \phi) &= \mu Y(\theta, \phi)
\end{align}
where $\lambda$ and $\mu$ are the eigenvalues to be determined.

\subsection{Eigenvalue Equation for $\hat{L}_z$}

We begin with the simpler operator, $\hat{L}_z$. In spherical coordinates, we derived in Section 4.2.1 that:
\begin{equation}
    \hat{L}_z = -i\hbar \frac{\partial}{\partial \phi}
\end{equation}

The eigenvalue equation is:
\begin{equation}
    -i\hbar \frac{\partial}{\partial \phi} Y(\theta, \phi) = \lambda Y(\theta, \phi)
\end{equation}

\subsubsection{Separation of Variables}
Since the operator only involves the variable $\phi$, we can separate the variables by proposing a solution of the form:
\begin{equation}
    Y(\theta, \phi) = f(\theta) \Phi(\phi)
\end{equation}

Substituting this into the differential equation:
\begin{equation}
    -i\hbar f(\theta) \frac{d\Phi(\phi)}{d\phi} = \lambda f(\theta) \Phi(\phi)
\end{equation}

Assuming $f(\theta) \neq 0$, we can divide by $f(\theta)$ to obtain a first-order ordinary differential equation for $\Phi(\phi)$:
\begin{equation}
    \frac{d\Phi}{d\phi} = \frac{i\lambda}{\hbar} \Phi(\phi)
\end{equation}

\subsubsection{General Solution}
The solution to this differential equation is an exponential function:
\begin{equation}
    \Phi(\phi) = A e^{i(\lambda/\hbar)\phi}
\end{equation}
where $A$ is a normalization constant. Consequently, the total eigenfunction takes the form:
\begin{equation}
    Y(\theta, \phi) = f(\theta) e^{i(\lambda/\hbar)\phi}
\end{equation}

\subsubsection{Boundary Conditions and Quantization}
For the wavefunction to be physically acceptable, it must be single-valued. In spherical coordinates, the angle $\phi$ and the angle $\phi + 2\pi$ represent the exact same physical point in space. Therefore, the function must satisfy the periodicity condition:
\begin{equation}
    \Phi(\phi + 2\pi) = \Phi(\phi)
\end{equation}

Substituting our solution:
\begin{equation}
    e^{i(\lambda/\hbar)(\phi + 2\pi)} = e^{i(\lambda/\hbar)\phi}
\end{equation}
\begin{equation}
    e^{i(\lambda/\hbar)\phi} e^{i(\lambda/\hbar)2\pi} = e^{i(\lambda/\hbar)\phi}
\end{equation}

Dividing both sides by $e^{i(\lambda/\hbar)\phi}$, we require:
\begin{equation}
    e^{2\pi i (\lambda/\hbar)} = 1
\end{equation}

From Euler's formula ($e^{ix} = \cos x + i\sin x$), this equality holds if and only if the exponent is an integer multiple of $2\pi i$. Let us define a number $m$ such that:
\begin{equation}
    \frac{\lambda}{\hbar} = m, \quad \text{where } m \in \mathbb{Z} = \{0, \pm 1, \pm 2, \dots\}
\end{equation}

\subsubsection{Result}
This boundary condition restricts the eigenvalues $\lambda$ to discrete values. We conclude that the eigenvalues of $\hat{L}_z$ are quantized in units of $\hbar$:
\begin{equation}
    \lambda = m\hbar
\end{equation}
And the angular dependence on $\phi$ is necessarily:
\begin{equation}
    Y(\theta, \phi) = f(\theta) e^{im\phi}
\end{equation}
where $m$ is called the \textit{magnetic quantum number}.

\subsection{Eigenvalue Equation for $\hat{L}^2$}

We now address the eigenvalue problem for the square of the angular momentum operator. Using the explicit expression for $\hat{L}^2$ in spherical coordinates derived in Eq. (62):
\begin{equation}
    \hat{L}^2 = -\hbar^2 \left[ \frac{1}{\sin\theta} \frac{\partial}{\partial \theta} \left( \sin\theta \frac{\partial}{\partial \theta} \right) + \frac{1}{\sin^2\theta} \frac{\partial^2}{\partial \phi^2} \right]
\end{equation}

We set up the eigenvalue equation aiming to find the eigenvalue $\lambda$:
\begin{equation}
    \hat{L}^2 Y(\theta, \phi) = \lambda Y(\theta, \phi)
\end{equation}

\subsubsection{Derivation of the Radial Equation}
We substitute the separable solution form obtained in the previous section, $Y(\theta, \phi) = f(\theta)e^{im\phi}$, into the eigenvalue problem.

First, we evaluate the action of the second derivative with respect to $\phi$ on our trial function:
\begin{equation}
    \frac{\partial^2}{\partial \phi^2} \left( f(\theta)e^{im\phi} \right) = f(\theta) \frac{\partial^2}{\partial \phi^2} e^{im\phi} = f(\theta) (im)^2 e^{im\phi} = -m^2 f(\theta)e^{im\phi}
\end{equation}

Substituting this result back into the full operator expression:
\begin{equation}
    -\hbar^2 \left[ \frac{1}{\sin\theta} \frac{\partial}{\partial \theta} \left( \sin\theta \frac{\partial f}{\partial \theta} \right) e^{im\phi} - \frac{m^2}{\sin^2\theta} f(\theta)e^{im\phi} \right] = \lambda f(\theta)e^{im\phi}
\end{equation}

Since the phase factor $e^{im\phi}$ is non-zero everywhere, we can divide the entire equation by it. Multiplying by $-1/\hbar^2$, we isolate the differential equation for the polar angle function $f(\theta)$:

\begin{equation}
    \frac{1}{\sin\theta} \frac{d}{d \theta} \left( \sin\theta \frac{d f}{d \theta} \right) - \frac{m^2}{\sin^2\theta} f(\theta) = -\frac{\lambda}{\hbar^2} f(\theta)
\end{equation}

Rearranging all terms to one side, we obtain the differential equation for $f(\theta)$:
\begin{equation}
    \frac{1}{\sin\theta} \frac{d}{d \theta} \left( \sin\theta \frac{d f}{d \theta} \right) + \left( \frac{\lambda}{\hbar^2} - \frac{m^2}{\sin^2\theta} \right) f(\theta) = 0
\end{equation}

\subsection{Natural Emergence of $\hat{L}^2$ and $\hat{L}_z$: Spherical Harmonics}

The differential equation derived above is the well-known \textit{Associated Legendre Differential Equation}. To see this clearly, one typically performs the change of variables $x = \cos\theta$, which maps the domain $\theta \in [0, \pi]$ to $x \in [-1, 1]$.

\subsubsection{Explicit Substitution}
Let us apply the chain rule to transform the derivatives from $\theta$ to $x$. Note that:
\begin{equation}
    \frac{d}{d\theta} = -\sin\theta \frac{d}{dx}
\end{equation}
Substituting this and the identity $\sin^2\theta = 1-x^2$ into the eigenvalue equation yields:
\begin{equation}
    \frac{d}{dx} \left[ (1-x^2) \frac{df}{dx} \right] + \left( \frac{\lambda}{\hbar^2} - \frac{m^2}{1-x^2} \right) f(x) = 0
\end{equation}

\subsubsection{Physical Admissibility: The Origin of Quantization}

The differential equation for $f(x)$ contains singular points at the poles $x = \pm 1$ (corresponding to $\theta = 0$ and $\theta = \pi$) due to the term $(1-x^2)$ in the denominator. 

Mathematically, the general solution to the Legendre equation is an infinite power series. Convergence analysis (specifically using the Frobenius method) reveals that this infinite series generally diverges logarithmically at the poles ($x \to \pm 1$). 

\paragraph{Regularity Condition}
For the wavefunction to represent a physical state, it must be single-valued, continuous, and finite everywhere (square-integrable). An infinite value at the poles is physically unacceptable. 
Therefore, we must impose a \textbf{regularity condition}: the infinite series must terminate to form a finite polynomial.

This truncation of the series only occurs if the coefficient of the recursion relation vanishes at some index $l$. This condition strictly demands that the eigenvalue parameter $\lambda$ satisfies:
\begin{equation}
    \frac{\lambda}{\hbar^2} = l(l+1), \quad l \in \mathbb{Z}^+
\end{equation}
Thus, the quantization of angular momentum magnitude ($\hat{L}^2$) is not an arbitrary postulate, but a direct consequence of the topological requirement that the wavefunction must be well-behaved on the closed surface of the sphere.

\paragraph{Constraints on $m$ via Rodrigues Formula}
With $\lambda$ fixed by $l$, the acceptable non-singular solutions are the Associated Legendre Polynomials $P_l^m(x)$. Their construction via the Rodrigues formula clearly elucidates the restriction on the magnetic quantum number $m$:

\begin{equation}
    P_l^m(x) \propto (1-x^2)^{m/2} \frac{d^{l+m}}{dx^{l+m}} \left[ (x^2-1)^l \right]
\end{equation}

Notice that the term $(x^2-1)^l$ is a polynomial of degree $2l$.
\begin{itemize}
    \item If we differentiate a polynomial of degree $2l$ more than $2l$ times, the result is zero.
    \item However, the structure of the Legendre equation specifically requires the order of the derivative to be $l+m$.
    \item If $|m| > l$, then $l+m > 2l$ (considering the structure of the differential operator), or more simply, the associated function involves derivatives that would annihilate the polynomial or produce singularities structure incompatible with the equation order.
\end{itemize}

Specifically, the Rodrigues definition implies that we can only take up to $l$ ``vertical'' steps (ladder operators) away from the base state. Therefore, for a non-vanishing solution, we must have:
\begin{equation}
    |m| \le l
\end{equation}
This explains why the projection of angular momentum ($\hat{L}_z$) cannot exceed its total magnitude ($\hat{L}^2$).

\subsubsection{Spherical Harmonics Definition}
Combining the radial part $P_l^m(\cos\theta)$ with the azimuthal phase $e^{im\phi}$, we obtain the \textbf{Spherical Harmonics} $Y_{l,m}(\theta, \phi)$ , which form the orthonormal basis on the sphere:

\begin{equation}
    Y_{l,m}(\theta, \phi) = (-1)^m \sqrt{\frac{(2l+1)}{4\pi} \frac{(l-m)!}{(l+m)!}} P_l^m(\cos\theta) e^{im\phi}
\end{equation}

\subsection{Simultaneous Eigenstates}

We conclude that there exist state vectors, denoted as $|l, m\rangle$, which are simultaneous eigenvectors of $\hat{L}^2$ and $\hat{L}_z$. These states completely describe the angular properties of the free particle (or any central potential system).

The defining eigenvalue equations are:
\begin{equation}
    \hat{L}^2 |l, m\rangle = \hbar^2 l(l+1) |l, m\rangle
\end{equation}
\begin{equation}
    \hat{L}_z |l, m\rangle = \hbar m |l, m\rangle
\end{equation}

These kets $|l, m\rangle$ form the angular momentum representation basis that we will use to solve the Schrödinger equation in the following section.